\begin{abstractpage}
	\begin{center}
		{\bf Rezumat}
	\end{center}
	\indent Lucrarea prezintă un sistem bazat pe Deep Learning pentru generarea automată a fișei tehnice a autoturismelor din gama Dacia 1300/1310, implementat cu PyTorch. Sistemul are o arhitectură pe două niveluri: primul nivel este un clasificator-router, care determină tipul imaginii (vedere frontală, vedere din spate sau compartimentul motor) și respinge imaginile pentru care încrederea este sub un prag configurabil. În funcție de rezultat, al doilea nivel interoghează automat clasificatorul specializat aferent: modelul (1300, 1310 CN3, CN4), tipul de caroserie (berlină sau break), respectiv tipul de alimentare (carburație sau injecție), iar predicțiile completează progresiv fișa tehnică a vehiculului. Toate clasificatoarele folosesc un backbone DINOv2 (Vision Transformer preantrenat self-supervised), complet înghețat, peste care se antrenează prin transfer learning câte un strat liniar, pe un set de date colectat și etichetat manual, în lipsa unui set de date public pentru aceste autoturisme. Performanța este măsurată prin validare încrucișată grupată pe vehicul, astfel încât evaluarea să reflecte comportamentul pe vehicule nevăzute la antrenare. Un astfel de sistem poate sprijini colecționarii și muzeele în identificarea configurației tehnice a vehiculelor istorice, etapă necesară, de exemplu, în evaluarea acestora.

	\begin{center}
		{\bf Abstract}
	\end{center}
	\indent This thesis presents a Deep Learning system, implemented in PyTorch, that automatically generates the technical specification sheet of cars from the Dacia 1300/1310 family. The system uses a two-level architecture: the first level is a router-classifier, which determines the type of the input image (front view, rear view, or engine bay) and rejects images whose confidence score falls below a configurable threshold. Depending on the result, the second level automatically queries the corresponding specialized classifier: the model variant (1300, 1310 CN3 and CN4), the body style (sedan or station wagon), or the fuel system (carburetor or fuel injection), and the predictions progressively fill in the vehicle's technical sheet. All classifiers share a frozen DINOv2 backbone (a self-supervised pretrained Vision Transformer), on top of which a single linear layer per task is trained via transfer learning, on a dataset that was collected and labeled manually, as no public dataset exists for these vehicles. Performance is measured with vehicle-grouped cross-validation, so that the evaluation reflects behaviour on vehicles unseen during training. Such a system can assist collectors and museums in identifying the technical configuration of historical vehicles, a necessary step, for instance, in their appraisal.
\end{abstractpage}
